\documentclass[aspectratio=169]{beamer}

\usetheme{Madrid}

\usepackage[T1]{fontenc}
\usepackage[utf8]{inputenc}
\usepackage{booktabs}
\usepackage{graphicx}
\usepackage{amsmath}
\usepackage{multicol}
\usepackage{ragged2e}

\definecolor{FormFixBlue}{RGB}{0,163,255}
\definecolor{DarkBG}{RGB}{14,17,23}
\definecolor{SoftGray}{RGB}{180,188,200}

\setbeamercolor{background canvas}{bg=DarkBG}
\setbeamercolor{normal text}{fg=white}
\setbeamercolor{frametitle}{bg=DarkBG, fg=FormFixBlue}
\setbeamercolor{title}{fg=FormFixBlue}
\setbeamercolor{progress bar}{fg=FormFixBlue, bg=gray}
\setbeamercolor{block title}{bg=FormFixBlue!18, fg=white}
\setbeamercolor{block body}{bg=white!4, fg=white}

\newcommand{\problemvisual}{\detokenize{/Users/syedahmed/Documents/Screenshot 2026-04-12 at 3.03.57 AM.png}}
\newcommand{\objectiveicon}{\detokenize{/Users/syedahmed/Documents/Screenshot 2026-04-12 at 3.06.01 PM.png}}
\newcommand{\overallarchitecture}{\detokenize{/Users/syedahmed/Documents/Screenshot 2026-04-12 at 3.05.49 PM.png}}
\newcommand{\trainingarchitecture}{\detokenize{/Users/syedahmed/Documents/Screenshot 2026-04-12 at 3.06.01 PM.png}}
\newcommand{\bilstmfig}{\detokenize{/Users/syedahmed/Downloads/Bidirectional LSTM-2026-04-12-094513.png}}
\newcommand{\lstmfig}{\detokenize{/Users/syedahmed/Downloads/Bidirectional LSTM-2026-04-12-101325.png}}
\newcommand{\transformerfig}{\detokenize{/Users/syedahmed/Downloads/Bidirectional LSTM-2026-04-12-101602.png}}
\newcommand{\accfig}{\detokenize{/Users/syedahmed/AI-Trainer/reports/analytics/accuracy_vs_epoch.png}}
\newcommand{\lossfig}{\detokenize{/Users/syedahmed/AI-Trainer/reports/analytics/loss_vs_epoch.png}}
\newcommand{\cmfig}{\detokenize{/Users/syedahmed/AI-Trainer/reports/analytics/confusion_matrix.png}}
\newcommand{\lstmtrainlog}{\detokenize{/Users/syedahmed/Documents/Screenshot 2026-04-13 at 12.29.23 AM.png}}
\newcommand{\transformertrainlog}{\detokenize{/Users/syedahmed/Documents/Screenshot 2026-04-13 at 12.26.58 AM.png}}
\newcommand{\workoutreport}{\detokenize{/Users/syedahmed/Documents/Screenshot 2026-04-12 at 3.03.57 AM.png}}

\title[FormFix AI]{FormFix AI}
\subtitle{Real-Time Exercise Pose Detection, Rep Counting, and Posture Correction}
\author[Ahmed \& Muradi]{Syed Ahmed \and Vibha V Muradi}
\institute[UVCE]{University of Visvesvaraya College of Engineering}
\date{May 4, 2026}

\begin{document}

\begin{frame}
    \titlepage
\end{frame}

\begin{frame}{The Challenge in Home Fitness}
    \begin{columns}[T]
        \begin{column}{0.52\textwidth}
            \textbf{The Supervision Gap}
            \begin{itemize}
                \item Improper form increases injury risk.
                \item Manual rep counting distracts the user.
                \item Most systems recognize the exercise but do not judge movement quality.
            \end{itemize}

            \vspace{0.3cm}
            \textbf{What FormFix AI Solves}
            \begin{itemize}
                \item real-time exercise classification,
                \item accurate rep counting,
                \item posture correction using joint angles,
                \item web-based progress tracking and workout reports.
            \end{itemize}
        \end{column}
        \begin{column}{0.44\textwidth}
            \centering
            \includegraphics[width=\textwidth,height=0.68\textheight,keepaspectratio]{\problemvisual}
        \end{column}
    \end{columns}
\end{frame}

\begin{frame}{Objectives: Beyond Recognition}
    \begin{columns}[T]
        \begin{column}{0.58\textwidth}
            \begin{enumerate}
                \item \textbf{Real-Time Perception:} detect skeleton landmarks from webcam video.
                \item \textbf{Temporal Intelligence:} compare LSTM, BiLSTM, and Transformer models.
                \item \textbf{Biomechanical Validation:} use joint-angle rules for posture correction.
                \item \textbf{Feedback Loop:} provide live voice and visual guidance.
                \item \textbf{Persistence:} store user workout progress and reports in SQL.
            \end{enumerate}
        \end{column}
        \begin{column}{0.38\textwidth}
            \centering
            \includegraphics[width=\textwidth,height=0.62\textheight,keepaspectratio]{\objectiveicon}
        \end{column}
    \end{columns}
\end{frame}

\begin{frame}{The Hybrid Intelligence Architecture}
    \centering
    \includegraphics[width=0.82\textwidth,height=0.76\textheight,keepaspectratio]{\overallarchitecture}
\end{frame}

\begin{frame}{From Video to Sequential Features}
    \begin{columns}[T]
        \begin{column}{0.48\textwidth}
            \textbf{Dataset Specs}
            \begin{itemize}
                \item Processed sequences: \textbf{1140}
                \item Features per frame: \textbf{132}
                \item Time window: \textbf{30 frames}
                \item Classes: \textbf{6}
            \end{itemize}

            \vspace{0.2cm}
            \textbf{Angle Formula}
            \[
            \theta = \cos^{-1}\left(\frac{(A-B)\cdot(C-B)}{\lVert A-B\rVert \, \lVert C-B\rVert}\right)
            \]

            \vspace{0.1cm}
            \textbf{Class Distribution}
            \begin{itemize}
                \item barbell\_squat: 152
                \item bicep\_curl: 225
                \item hammer\_curl: 78
                \item lateral\_raise: 170
                \item pull\_up: 212
                \item push\_up: 303
            \end{itemize}
        \end{column}
        \begin{column}{0.48\textwidth}
            \centering
            \includegraphics[width=\textwidth,height=0.72\textheight,keepaspectratio]{\trainingarchitecture}
        \end{column}
    \end{columns}
\end{frame}

\begin{frame}{Model Strategy: Why Three Models?}
    \begin{table}
        \centering
        \small
        \begin{tabular}{p{2.2cm} p{3.2cm} p{5.1cm}}
            \toprule
            \textbf{Model} & \textbf{Role} & \textbf{Reason for Use} \\
            \midrule
            LSTM & Baseline model & Lightweight and fast sequence model for comparison. \\
            BiLSTM & Main temporal model & Uses forward and backward context; stronger motion understanding. \\
            Transformer & Advanced comparison & Uses self-attention for long-range frame relationships. \\
            \bottomrule
        \end{tabular}
    \end{table}

    \vspace{0.3cm}
    \centering
    \textbf{Deployment logic:} the system supports \textit{Fast}, \textit{Best}, \textit{Hybrid}, and \textit{Ensemble} modes.\\
    In practice, BiLSTM is weighted highest in the hybrid pipeline.
\end{frame}

\begin{frame}{LSTM, BiLSTM and Transformer Architectures}
    \begin{columns}[T]
        \begin{column}{0.32\textwidth}
            
\includegraphics[width=\textwidth,height=0.24\textheight,keepaspectratio]{\lstmfig}
            \vspace{0.08cm}
            \scriptsize \textbf{LSTM:} single-direction temporal baseline.
        \end{column}
        \begin{column}{0.32\textwidth}
            
\includegraphics[width=\textwidth,height=0.24\textheight,keepaspectratio]{\bilstmfig}
            \vspace{0.08cm}
            \scriptsize \textbf{BiLSTM:} bidirectional temporal classifier.
        \end{column}
        \begin{column}{0.32\textwidth}
            
\includegraphics[width=\textwidth,height=0.24\textheight,keepaspectratio]{\transformerfig}
            \vspace{0.08cm}
            \scriptsize \textbf{Transformer:} attention-based sequence model.
        \end{column}
    \end{columns}

    \vspace{0.25cm}
    \begin{block}{Working Principle}
    All three models receive the same 30-frame pose-angle sequence, but they learn motion dependencies differently. They are trained independently and then compared for speed, stability, and accuracy.
    \end{block}
\end{frame}

\begin{frame}{Validation Engines: Counting and Correction}
    \begin{multicols}{2}
        \textbf{Repetition Counting}
        \begin{itemize}
            \item exercise-specific tracked angle,
            \item smoothing and direction gating,
            \item stage transitions such as up/down,
            \item threshold validation before incrementing rep count.
        \end{itemize}

        \columnbreak

        \textbf{Posture Correction}
        \begin{itemize}
            \item elbow, knee, shoulder, and hip angles,
            \item rule templates for each exercise,
            \item mistake codes such as elbow instability or insufficient depth,
            \item live voice cue and end-of-exercise report.
        \end{itemize}
    \end{multicols}

    \vspace{0.2cm}
    \centering
    \includegraphics[width=0.78\textwidth,height=0.3\textheight,keepaspectratio]{\problemvisual}
\end{frame}

\begin{frame}{Training and Evaluation Snapshots}
    \begin{columns}[T]
        \begin{column}{0.49\textwidth}
            \includegraphics[width=\textwidth,height=0.72\textheight,keepaspectratio]{\lstmtrainlog}
        \end{column}
        \begin{column}{0.49\textwidth}
            \includegraphics[width=\textwidth,height=0.72\textheight,keepaspectratio]{\transformertrainlog}
        \end{column}
    \end{columns}
\end{frame}

\begin{frame}{Quantitative Results}
    \begin{columns}[T]
        \begin{column}{0.42\textwidth}
            \textbf{Current Results}
            \begin{itemize}
                \item LSTM Top-1: \textbf{84.30\%}
                \item BiLSTM Top-1: \textbf{93.72\%}
                \item BiLSTM Top-3: \textbf{99.82\%}
                \item Transformer Top-1: \textbf{97.89\%}
            \end{itemize}

            \vspace{0.2cm}
            \textbf{BiLSTM Quality Metrics}
            \begin{itemize}
                \item Weighted F1-score: \textbf{93.64\%}
                \item Macro F1-score: \textbf{91.69\%}
            \end{itemize}

            \vspace{0.2cm}
            \textbf{Strongest classes}
            \begin{itemize}
                \item push\_up: F1 \textbf{0.997}
                \item pull\_up: recall \textbf{1.00}
                \item barbell\_squat: F1 \textbf{0.94}
            \end{itemize}
        \end{column}
        \begin{column}{0.54\textwidth}
            \centering
            \includegraphics[width=\textwidth,height=0.72\textheight,keepaspectratio]{\cmfig}
        \end{column}
    \end{columns}
\end{frame}

\begin{frame}{How Accuracy Is Calculated}
    \small
    \textbf{Top-1 Accuracy}
    \[
    \text{Top-1} = \frac{\text{Correct highest-confidence predictions}}{\text{Total predictions}} \times 100
    \]

    \textbf{Top-3 Accuracy}
    \[
    \text{Top-3} = \frac{\text{True class inside top 3 predictions}}{\text{Total predictions}} \times 100
    \]

    \textbf{Precision, Recall, F1-score}
    \[
    \text{Precision}=\frac{TP}{TP+FP},\quad
    \text{Recall}=\frac{TP}{TP+FN},\quad
    \text{F1}=2\cdot\frac{PR}{P+R}
    \]

    \vspace{0.2cm}
    \begin{columns}[T]
        \begin{column}{0.49\textwidth}
            \includegraphics[width=\textwidth,height=0.28\textheight,keepaspectratio]{\accfig}
        \end{column}
        \begin{column}{0.49\textwidth}
            \includegraphics[width=\textwidth,height=0.28\textheight,keepaspectratio]{\lossfig}
        \end{column}
    \end{columns}
\end{frame}

\begin{frame}{Integrated Web Application}
    \begin{columns}[T]
        \begin{column}{0.46\textwidth}
            \centering
            \includegraphics[width=\textwidth,height=0.68\textheight,keepaspectratio]{\workoutreport}
        \end{column}
        \begin{column}{0.5\textwidth}
            \textbf{Implemented System Components}
            \begin{itemize}
                \item FastAPI backend for live frame analysis,
                \item React frontend for workout and nutrition flows,
                \item SQL-based user progress tracking,
                \item meal planner, scan-to-calorie flow, and recipe CRUD,
                \item session reports with perfect reps, corrected reps, and common mistakes.
            \end{itemize}

            \vspace{0.2cm}
            \textbf{Improvement over surveyed papers}
            \begin{itemize}
                \item not only recognition,
                \item not only pose correction,
                \item but a full usable coaching application.
            \end{itemize}
        \end{column}
    \end{columns}
\end{frame}

\begin{frame}{Final Summary and Future Work}
    \begin{block}{Key Takeaway}
        FormFix AI turns research ideas from pose estimation, sequential learning, and biomechanical validation into a practical real-time fitness trainer.
    \end{block}

    \begin{columns}[T]
        \begin{column}{0.5\textwidth}
            \textbf{Future Improvements}
            \begin{itemize}
                \item expand to deadlift and overhead press,
                \item create a stricter unseen holdout benchmark,
                \item improve mobile deployment,
                \item add multi-view support,
                \item improve browser-side speed and user-specific calibration.
            \end{itemize}
        \end{column}
        \begin{column}{0.45\textwidth}
            \textbf{Conclusion}
            \begin{itemize}
                \item strong real-time pipeline,
                \item high current classification performance,
                \item deployable UI plus database integration,
                \item clear scope for improvement and scale.
            \end{itemize}
        \end{column}
    \end{columns}
\end{frame}

\begin{frame}
    \centering
    \Huge \textcolor{FormFixBlue}{Thank You}

    \vspace{0.4cm}
    \Large \textbf{FormFix AI}

    \vspace{0.2cm}
    \small Questions?
\end{frame}

\end{document}
